% Mit \section{...} eröffnen wir einen neuen Abschnitt.
% Der Befehl setzt nicht nur den Text in einer größeren,
% fetten Schrift, sondern sorgt außerdem dafür, daß er im
% Inhaltsverzeichnis erscheint.
%
% Mit \label{...} erzeugen wir einen Bezeichner, mit dessen Hilfe
% wir später auf die Nummer des Abschnitts verweisen können (nämlich
% mit \ref{...}).
%
% Das Kommentarzeichen hinter „Übersicht“ dient dazu, ein
% Leerzeichen zwischen „Übersicht“ und dem \label-Befehl
% zu vermeiden, das andernfalls sichtbar würde – z.B. im
% Inhaltsverzeichnis.
%
\section{Einleitende Übersicht%
         \label{sec:Einleitung}}

Diese Vorlage ist dazu gedacht,
den Einstieg in das Schreiben von größeren wissenschaftlichen Dokumenten
mit Hilfe von \LaTeX{} zu erleichtern,
insbesondere z.\,B.~für eine Bachelorarbeit
  % Zwischen einbuchstabigen Abkürzungen verwendet man halbe Leerzeichen,
  % die man in LaTeX mit dem Befehl „\,“ (Backslash und Komma) erzeugt.
  % Der Punkt nach der Abkürzung soll nicht mit einem Punkt am Ende
  % des Satzes verwechselt werden, daher verwenden wir hierfür entweder 
  % „\ “ (Backslash und Leerzeichen) oder „~“ (Tilde). Ersteres erlaubt
  % einen Zeilenumbruch an diesem Leerzeichen, letzteres nicht. In diesem 
  % konkreten Beispiel bewirkt also die Tilde, daß das „z.B.“ nicht am
  % Ende der Zeile stehen darf.
am Campus Velbert-Heiligenhaus (CVH) der Hochschule Bochum (HS Bochum). 
Die erste Version ist durch Jan Weber erarbeitet worden.
Ihm gebührt der Dank für den größten Arbeitsanteil an dieser Dokumentationsvorlage.
Die Überarbeitung der ersten Version erfolgte durch Markus Lemmen und Peter Gerwinski. 

Um ein derartiges Ziel erreichen zu können,
liefert die vorliegende Arbeit de facto eine Vorlage
für beliebige am CVH zu erstellende wissenschaftliche Ausarbeitungen. 
Für andere Arbeiten als die hier ausschließlich ausgeführte
Vorlage einer Bachelorarbeit sind verschiedene Details abzuändern
wie z.\,B.\ das Deckblatt. 
Die eigentliche Struktur einer Arbeit kann jedoch bestehen bleiben,
wie auch die grundlegende Struktur einer wissenschaftlichen Ausarbeitung in 

\begin{enumerate}
  \item Einleitende Übersicht, 
  \item Grundlagen, 
  \item Konzeption, 
  \item Prototypische Implementierung mit eventueller Diskussion sowie
  \item Zusammenfassung und Ausblick. 
\end{enumerate}

Neben der Bereitstellung einer derartigen Strukturierung von
wissenschaftlichen Arbeiten verfolgt die vorliegende Dokumentation
auch das Ziel, Beispiele für den Gebrauch und die Anwendung von
\LaTeX\ innerhalb dieser Struktur zu geben und so die Nutzung des
Textsatzsystems \LaTeX\ zu erleichtern und entsprechende Hürden
abzubauen. 
Eine umfassende Einleitung in \LaTeX\ oder das wissenschaftliche
Schreiben kann diese Vorlage jedoch nicht geben. 
In diesem Zusammenhang sei auf die vielen Quellen im Internet,
Büchern und auch auf die entsprechenden Kurse am Campus verwiesen. 

Dementsprechend gliedert sich die vorliegende Arbeit wie folgt: 
Im Anschluss an diese einleitende Übersicht im
Abschnitt~\ref{sec:Einleitung} folgt die Darstellung der wichtigsten
  % \ref{...}: Hier verwenden wir den oben definierten \label{...}.
\LaTeX\ Befehle und Umgebungen in Abschnitt \ref{sec:LaTeX-Beispiele}
  % Dieser Label wird erst im nächsten Abschnitt definiert.
  % LaTeX löst diese Vorwärts-Definition durch das Anlegen
  % temporärer Dateien auf, die beim nächsten Compilieren
  % eingebunden werden. Dies ist der Grund, weshalb Referenzen
  % gelegentlich erst nach mehrfachem Compilieren korrekt sind.
mit Hilfe vieler konkreter Beispielen.
Nach einer knappen Zusammenfassung des Compilierens in Abschnitt~\ref{sec:Compilieren}
und einem Exkurs in die Strukturierung in Abschnitte in Abschnitt~\ref{sec:Abschnitte}
gibt Abschnitt~\ref{sec:Bilder-Tabellen} einen ersten Einblick,
wie mit Hilfe von Umgebungen Bilder und Tabellen generell
bzw.\ in Abschnitt~\ref{sec:Bilder}
  % Auch der Punkt hinter der Abkürzung „bzw.“ ist kein Punkt am
  % Ende des Satzes – daher das „\ “ (Backslash und Leerzeichen).
  % Die Tilde „~“ wirkt genauso, verbietet aber einen Zeilenumbruch
  % an dieser Stelle, daher steht sie zwischen „Abschnitt“ und der
  % Referenz: Wir wollen nicht, daß „Abschnitt“ als letztes Wort in
  % der Zeile steht und die Referenz (eine Zahl) in der nächsten
  % Zeile.
spezielle Anordnungen von Bildern anzuordnen sind.
Abschnitt~\ref{sec:Tabellen} erklärt Tabellen,
Abschnitt~\ref{sec:Gleichungen} zeigt Gleichungen,
und Abschnitt~\ref{sec:Nomenklatur} veranschaulicht Einträge in die Nomenklatur,
häufig auch „Formelzeichen und Bezeichnungen“ genannt.
Eine Zusammenfassung und einen Ausblick schließen diese Arbeit in
Abschnitt~\ref{sec:zusammen} ab. Eventuell noch benötigte Anhänge
finden sich in den Anhängen \ref{anh:fertigungszeichnungen} bis
\ref{anh:datenblaetter}. 
